\subsection{GRADED COMMUTATIVE GROUPS}

We are going to translate into another language the definitions concerning direct sums (§ 1, no. 8).

DEFINITION 1. Given a commutative group G written additively and a set \(\Delta\) , a graduation of type \(\Delta\) on G is a family \((\mathbf{G}_{\lambda})_{\lambda \in L}\) of subgroups of G, of which G is the direct sum. The set G, with the structure defined by its group law and its graduation, is called a graded (commutative) group of type \(\Delta\) .

\(\Delta\) is called the set of degrees of \(\mathbf{G}\) . An element \(x \in \mathbf{G}\) is called homogeneous if it belongs to one of the \(\mathbf{G}_{\lambda}\) , homogeneous of degree \(\lambda\) if \(x \in \mathbf{G}_{\lambda}\) . The element 0 is therefore homogeneous of all degrees; but if \(x \neq 0\) is homogeneous, it belongs to only one of the \(\mathbf{G}_{\lambda}\) ; the index \(\lambda\) such that \(x \in \mathbf{G}_{\lambda}\) is then called the degree of \(x\) (or sometimes the weight of \(x\) ) and is sometimes denoted by \(\deg(x)\) . Every \(y \in \mathbf{G}\)

may be written uniquely as a sum \(\sum_{\lambda} y_{\lambda}\) of homogeneous elements with \(y_{\lambda} \in G_{\lambda}\) ; \(y_{\lambda}\) is called the homogeneous component of degree \(\lambda\) (or simply the component of degree \(\lambda\) ) of \(y\) . When the word ``weight'' is used instead of ``degree'', the adjective ``homogeneous'' is replaced by ``isobaric''.

Examples. (1) Given any commutative monoid \(\Delta\) (with identity element 0) and a commutative group \(G\) , a graduation \((G_{\lambda})_{\lambda \in \Delta}\) is defined on \(G\) by taking \(G_0 = G\) and \(G_{\lambda} = \{0\}\) for \(\lambda \neq 0\) ; this graduation is called trivial.

\begin{enumerate}
\def\labelenumi{(\arabic{enumi})}
\setcounter{enumi}{1}
\tightlist
\item
  Let \(\Delta, \Delta'\) be two sets and \(\rho\) a mapping of \(\Delta\) into \(\Delta'\) . Let \((\mathbf{G}_{\lambda})_{\lambda \in \Delta}\) be a graduation of type \(\Delta\) on a commutative group \(\mathbf{G}\) ; for \(\mu \in \Delta'\) , let \(\mathbf{G}_{\mu}'\) be the sum of the \(\mathbf{G}_{\lambda}\) such that \(\rho(\lambda) = \mu\) ; clearly \((\mathbf{G}_{\mu}')_{\mu \in \Delta'}\) is a graduation of type \(\Delta'\) on \(\mathbf{G}\) , said to be derived from \((\mathbf{G}_{\lambda})\) by means of the mapping \(\rho\) .
\end{enumerate}

When \(\Delta\) is a commutative group written additively and \(\rho\) the mapping \(\lambda \mapsto -\lambda\) of \(\Delta\) onto itself, \((\mathbf{G}_{\mu}^{\prime})\) is called the opposite graduation of \((\mathbf{G}_{\lambda})\) .

\begin{enumerate}
\def\labelenumi{(\arabic{enumi})}
\setcounter{enumi}{2}
\item
  If \(\Delta = \Delta_1 \times \Delta_2\) is a product of two sets, a graduation of type \(\Delta\) is called a bigraduation of types \(\Delta_1, \Delta_2\) . For all \(\lambda \in \Delta_1\) , let \(G_{\lambda}' = \bigoplus_{\mu \in \Delta_2} G_{\lambda\mu}\) and, for all \(\mu \in \Delta_2\) , let \(G_{\mu}' = \bigoplus_{\lambda \in \Delta_1} G_{\lambda\mu}\) ; clearly \((G_{\lambda}')_{\lambda \in \Delta_1}\) is a graduation of type \(\Delta_1\) and \((G_{\mu}'')_{\mu \in \Delta_2}\) a graduation of type \(\Delta_2\) on \(G\) ; these graduations are called the partial graduations derived from the bigraduation \((G_{\lambda\mu})\) . Note that \(G_{\lambda\mu} = G_{\lambda}' \cap G_{\mu}'\) ; conversely, if \((G_{\lambda}')_{\lambda \in \Delta_1}\) and \((G_{\mu}'')_{\mu \in \Delta_2}\) are two graduations on \(G\) such that \(G\) is the direct sum of the \(G_{\lambda\mu} = G_{\lambda}' \cap G_{\mu}'\) , these subgroups form a bigraduation of types \(\Delta_1, \Delta_2\) on \(G\) , of which \((G_{\lambda}')\) and \((G_{\mu}'')\) are the partial graduations. We leave to the reader the task of generalizing this to the case where \(\Delta\) is a finite product of sets.
\item
  Let \(\Delta_0\) be a commutative monoid written additively, with identity element denoted by 0; let I be any set and \(\Delta_0^{(I)} = \Delta\) denote the submonoid of the product \(\Delta_0^I\) consisting of the families \((\lambda_t)_{t \in I}\) of finite support. Let \(\rho: \Delta \to \Delta_0\) be the surjective (codiagonal) homomorphism of \(\Delta\) into \(\Delta_0\) defined by \(\rho((\lambda_t)) = \sum_{t \in I} \lambda_t\) . From every graduation of type \(\Delta\) a graduation of type \(\Delta_0\) is derived by means of \(\rho\) (Example 2); it is called the total graduation associated with the given ``multigraduation'' of type \(\Delta\) .
\end{enumerate}

The definitions and examples of this no. extend immediately to the case where G is a group which is not necessarily commutative; it is simply necessary to replace everywhere the notion of direct sum by that of ``restricted sum'' (§ 1, no. 6, Remark). Note that in this case the \(G_{\lambda}\) are normal subgroups of G and that for \(\lambda \neq \mu\) every element of \(G_{\lambda}\) is permutable with every element of \(G_{\mu}\) .

